%! Fundamental Theorem of Algebra and Complex Zeros — Unit 3, Lesson 3.5 — Slides (Beamer)
\documentclass[aspectratio=169, 11pt]{beamer}
\usepackage{algebra2-beamer}   % provides \forestheader{} and \sectionlabel[color]{}

\begin{document}

% TITLE SLIDE (CourseName is not defined in beamer, so the name is literal)
{
\setbeamercolor{background canvas}{bg=forest}
\begin{frame}[plain]
  \vspace{2.8cm}\hspace{0.55cm}%
  \begin{minipage}{0.9\textwidth}
    {\color{goldacc}\bfseries\small LESSON 3.5}\\[0.5cm]
    {\color{white}\bfseries\LARGE Fundamental Theorem of Algebra}\\[0.25cm]
    {\color{white}\bfseries\large and Complex Zeros}\\[1.0cm]
    {\color{forestmid}\small Algebra 2~$\cdot$~Unit 3: Polynomial Functions}
  \end{minipage}
\end{frame}
}

% HOOK
\begin{frame}[t, plain]
  \forestheader{One graph, one intercept, three solutions}
  \sectionlabel{Hook}
  \begin{columns}[T]
    \begin{column}{0.52\textwidth}
      \[ f(x) = x^3 - x^2 + 4x - 4 \]
      \vspace{2pt}
      \begin{itemize}\setlength{\itemsep}{4pt}
        \item Yesterday it stalled: $x=1$ worked, and the depressed polynomial was $x^2+4$.
        \item In the Warm-Up you finished it: $x = 1$, $2i$, $-2i$.
        \item But its graph crosses the $x$-axis \textbf{exactly once}.
      \end{itemize}
    \end{column}
    \begin{column}{0.46\textwidth}
      \begin{block}{The question}
        A cubic is supposed to have three solutions. The graph shows one. \\[4pt]
        \textbf{Where are the other two?}
      \end{block}
      \vspace{3pt}
      {\footnotesize All year, ``no solution'' has really meant \textbf{no \emph{real} solution.}
       Today that sentence gets finished.}
    \end{column}
  \end{columns}
\end{frame}

% THE THEOREM
\begin{frame}[t, plain]
  \forestheader{The Fundamental Theorem of Algebra}
  \sectionlabel[goldacc]{Foundation}
  \begin{block}{The theorem}
    Every polynomial of degree $n \ge 1$ has \textbf{at least one} complex zero.
  \end{block}
  \vspace{2pt}
  \begin{block}{The corollary we actually use}
    A polynomial of degree $n$ has \textbf{exactly $n$} complex zeros, counted with
    \textbf{multiplicity}, and factors completely as
    \[ f(x) = a(x-r_1)(x-r_2)\cdots(x-r_n). \]
  \end{block}
  \vspace{2pt}
  {\footnotesize \textbf{Why:} find a zero, divide it out (Factor Theorem, 3.3), and the quotient has
   degree $n-1$. Repeat. The process runs exactly $n$ times.}
\end{frame}

% COUNTING WITH MULTIPLICITY
\begin{frame}[t, plain]
  \forestheader{``Counted with multiplicity'' is not fine print}
  \sectionlabel{Counting}
  \begin{columns}[T]
    \begin{column}{0.5\textwidth}
      \footnotesize
      The unit anchor, from Lesson 3.0:
      \[ g(x) = x^3 - 3x + 2 = (x-1)^2(x+2) \]
      \vspace{-6pt}
      \begin{itemize}\setlength{\itemsep}{5pt}
        \item \textbf{Distinct} zeros: $1$ and $-2$ --- only two.
        \item \textbf{With multiplicity}: $1$, $1$, $-2$ --- exactly three. \checkmark
        \item The graph \emph{touches} at $x=1$ (even multiplicity) and \emph{crosses} at $x=-2$.
      \end{itemize}
    \end{column}
    \begin{column}{0.48\textwidth}
      \begin{block}{The sentence to remember}
        Distinct zeros can be fewer than $n$. \\[4pt]
        Zeros counted with multiplicity are never fewer and never more --- always exactly $n$.
      \end{block}
    \end{column}
  \end{columns}
\end{frame}

% CONJUGATE PAIRS
\begin{frame}[t, plain]
  \forestheader{Imaginary zeros travel in pairs}
  \sectionlabel[goldacc]{Theorem}
  \begin{block}{Complex Conjugate Root Theorem}
    If a polynomial has \textbf{real} coefficients and $a+bi$ is a zero ($b \ne 0$), then $a-bi$ is a
    zero too.
  \end{block}
  \vspace{3pt}
  \begin{columns}[T]
    \begin{column}{0.52\textwidth}
      \footnotesize
      \textbf{Why the pair is forced:}
      \[ \big((x-a)-bi\big)\big((x-a)+bi\big) = (x-a)^2 + b^2 \]
      \[ = x^2 - 2ax + (a^2+b^2) \]
      Every $i$ cancels. A \emph{lone} imaginary zero would strand an $i$ in the coefficients.
    \end{column}
    \begin{column}{0.44\textwidth}
      \footnotesize
      \textbf{Warm-Up evidence:}\\[3pt]
      $x^2-2x+5 \Rightarrow 1+2i$ and $1-2i$ \\[3pt]
      $x^2+4 \Rightarrow 2i$ and $-2i$ \\[6pt]
      Not a coincidence --- a requirement.
    \end{column}
  \end{columns}
\end{frame}

% CONSEQUENCES
\begin{frame}[t, plain]
  \forestheader{Two consequences worth more than the theorem}
  \sectionlabel{Reasoning}
  \begin{itemize}\setlength{\itemsep}{8pt}
    \item The number of imaginary zeros is always \textbf{even} --- $0$, $2$, $4$, \dots --- because
      they are counted in pairs.
    \item Therefore a polynomial of \textbf{odd} degree with real coefficients has \textbf{at least
      one real zero}. Odd degree minus an even number of imaginary zeros leaves an odd number of real
      ones, and that is never $0$.
  \end{itemize}
  \vspace{4pt}
  \begin{block}{Lesson 3.0 said the same thing with a picture}
    An odd-degree graph's two arms point in \emph{opposite} directions --- so the graph has no choice
    but to cross the $x$-axis. Two arguments, one conclusion.
  \end{block}
\end{frame}

% NUMBER AND TYPE
\begin{frame}[t, plain]
  \forestheader{Number and type --- without solving anything}
  \sectionlabel[goldacc]{The Standard}
  \[ \text{imaginary zeros} \;=\; \text{degree} \;-\; \text{real zeros (counted with multiplicity)} \]
  \vspace{-2pt}
  \begin{columns}[T]
    \begin{column}{0.52\textwidth}
      \footnotesize
      \begin{itemize}\setlength{\itemsep}{5pt}
        \item An $x$-intercept \emph{is} a real zero.
        \item An imaginary zero is \textbf{nowhere} on the graph.
        \item A cubic crossing once: $3-1 = \mathbf{2}$ imaginary.
        \item A quartic crossing twice: $4-2 = \mathbf{2}$ imaginary.
        \item A quintic crossing once: $5-1 = \mathbf{4}$ imaginary.
      \end{itemize}
    \end{column}
    \begin{column}{0.44\textwidth}
      \begin{block}{Check yourself}
        Could a \emph{cubic} with real coefficients have \textbf{no} $x$-intercepts? \\[4pt]
        No --- that needs three imaginary zeros, and three is odd.
      \end{block}
    \end{column}
  \end{columns}
\end{frame}

% SOLVING OVER C
\begin{frame}[t, plain]
  \forestheader{The same workflow --- with a different last step}
  \sectionlabel{Method}
  \begin{columns}[T]
    \begin{column}{0.5\textwidth}
      \footnotesize
      \textbf{Solve} $x^3-5x^2+9x-5=0$. \\[4pt]
      \textbf{1--3. List, test, divide:} $r=1$ works.
      \begin{center}
      \renewcommand{\arraystretch}{1.3}
      \begin{tabular}{r|rrrr}
      $1$ & $1$ & $-5$ & $9$ & $-5$ \\
          &     & $1$  & $-4$ & $5$ \\ \cline{2-5}
          & $1$ & $-4$ & $5$  & $0$ \\
      \end{tabular}
      \end{center}
      \textbf{4. Finish:} $x^2-4x+5$, $D = -4$, \\
      so $x = \dfrac{4 \pm 2i}{2} = 2 \pm i$.
    \end{column}
    \begin{column}{0.48\textwidth}
      \begin{block}{Answer}
        $x = 1,\; 2+i,\; 2-i$ \\[4pt]
        One real, two imaginary. $1+2 = 3 =$ the degree. \checkmark
      \end{block}
      \vspace{3pt}
      {\footnotesize The only thing that changed: a negative discriminant is no longer ``no
       solution.'' It is a \textbf{conjugate pair}. \\[4pt]
       And $x^4-5x^2-36 = (x^2-9)(x^2+4)$ needs no division at all: $\pm3, \pm2i$.}
    \end{column}
  \end{columns}
\end{frame}

% BUILDING BACKWARD
\begin{frame}[t, plain]
  \forestheader{Run it backward: zeros $\rightarrow$ factors $\rightarrow$ polynomial}
  \sectionlabel[goldacc]{Building}
  \begin{columns}[T]
    \begin{column}{0.5\textwidth}
      \footnotesize
      \textbf{(a)} Zeros $-1$ (multiplicity 2) and $4$:
      \[ (x+1)^2(x-4) = x^3-2x^2-7x-4 \]
      \vspace{-8pt}
      \textbf{(b)} Real coefficients, zeros $2$ and $3i$: \\[3pt]
      $-3i$ has to come too, and
      \[ (x-2)(x^2+9) = x^3-2x^2+9x-18 \]
    \end{column}
    \begin{column}{0.48\textwidth}
      \begin{block}{Each zero $r$ gives the factor $(x-r)$}
        A conjugate pair gives the \emph{real} quadratic $x^2-2ax+(a^2+b^2)$ --- multiply the pair
        together first and the $i$'s never appear.
      \end{block}
      \vspace{3pt}
      {\footnotesize \textbf{Test question:} smallest degree of a real polynomial with zeros $5$ and
       $1-i$? \; \textbf{3} --- because $1+i$ is not optional.}
    \end{column}
  \end{columns}
\end{frame}

% WHERE THIS GOES
\begin{frame}[t, plain]
  \forestheader{Nothing was ever missing}
  \sectionlabel{Connections}
  \begin{columns}[T]
    \begin{column}{0.5\textwidth}
      \footnotesize
      \textbf{Today:}
      \begin{itemize}\setlength{\itemsep}{3pt}
        \item Degree $n$ $\Rightarrow$ exactly $n$ zeros, with multiplicity.
        \item Real coefficients $\Rightarrow$ imaginary zeros in conjugate pairs.
        \item Number and type, from a degree and a graph.
        \item Solve over $\mathbb{C}$; build from a list of zeros.
      \end{itemize}
    \end{column}
    \begin{column}{0.48\textwidth}
      \begin{block}{Coming next}
        \textbf{3.6} Graphing polynomial functions \& modeling --- end behavior (3.0), zeros
        (3.2--3.4), multiplicity (3.0), and today's complete count become \emph{one} procedure for
        producing and reading a polynomial's graph.
      \end{block}
    \end{column}
  \end{columns}
\end{frame}

\end{document}
